\section{Introduction}
\label{sect:intro}

Binary-star orbits are among the principal routes to dynamical stellar masses. Independent resolved relative astrometry constrains the sky-plane geometry, while a double-lined spectroscopic binary (SB2) supplies both line-of-sight velocity amplitudes and strong information on the mass ratio. Gaia adds precise one-dimensional astrometric information over a mission-length time baseline \citep{GaiaMission2016}. The scientific issue addressed here is not how to combine those data types in general, but how the \emph{measurement model} used for a luminous close pair propagates into the inferred physical orbit, parallax, light ratio, and component masses.

The broad combined-orbit problem is already well developed. BINARYS jointly treats Hipparcos/Gaia absolute astrometry, independent relative astrometry, and component-labelled radial velocities \citep{Leclerc2023}, while orvara performs posterior inference from radial velocity, absolute astrometry, and relative astrometry \citep{Brandt2021}. \citet{Chevalier2023} used BINARYS to combine Gaia DR3 non-single-star solutions with SB2 data and derive individual masses and luminosities, demonstrating directly that Gaia astrometry plus SB2 spectroscopy is not itself a gap. Component-mass inference from unresolved Gaia astrometric orbits and photometric information is also established \citep{BailerJones2026}.

The unresolved-to-resolved transition remains more problematic. Gaia's published DR3 astrometric-binary processing used a preliminary treatment to discard partially resolved double stars before orbital modelling \citep{Halbwachs2023}, and \citet{Holl2023} demonstrated that close source structure can introduce scan-angle-dependent image-parameter biases and spurious signals. BINARYS is particularly informative in this context: it already models non-trivial Hipparcos transit response when secondary light matters, but its Gaia treatment does not extend the same response modelling to luminous partially resolved Gaia pairs; \citet{Leclerc2023} explicitly identify the need for more detailed Gaia information to treat that flux-contamination regime.

A nonlinear separation-dependent Gaia-like blended-source response is also prior art. \citet{ElBadry2024} released \texttt{gaiamock}, whose \texttt{al\_bias\_binary} places the measured one-dimensional coordinate at the peak of the combined along-scan light profile. Its public implementation also contains response-aware stellar astrometry-plus-RV forward prediction. We therefore make no claim that resolution-aware stellar Gaia modelling, or response-aware Gaia-plus-RV prediction, originates here. Our equal-width implementation reproduces the same response family over the common single-peak domain.

Nor is response-aware orbital inference itself new in a broader astronomical sense. \citet{Liu2024} fitted the binary asteroid (4337) Arecibo using Gaia Focused Product Release astrometry with an unresolved/partially-resolved/resolved response and external occultation constraints. Their analysis is not a stellar SB2 joint inference, but it removes any defensible claim that putting a partially resolved Gaia response inside orbital inference is new in general.

The response theory is also more general than a one-dimensional equal-width model. \citet{Penoyre2026} derived positions and resolvability for blended Gaussian point sources observed with an elongated Gaia-like PSF and showed that the problem reduces to a one-dimensional blend after scaling by an orientation-dependent effective width. The published correction to that work changes signs in several low-order expressions; the implementation used here therefore employs the exact effective-width relation and numerical peak solution rather than those series approximations. Gaia itself uses a substantially richer calibrated point/line-spread function with drift-scan, colour, and focal-plane dependences \citep{Rowell2021,Rowell2026}. None of the Gaussian widths in this work is identified with a calibrated Gaia resolution scale.

Bayesian Gaia epoch-astrometry plus radial-velocity inference is likewise already established in \texttt{kima} \citep{Baycroft2026}. Against this background, we investigate a deliberately narrow \emph{candidate inference gap}. Based on the published literature and public implementations examined for this project, we did not identify a demonstrated stellar target-level analysis in which a physical marginal-resolution Gaia response is placed directly inside an epoch-astrometric likelihood while the same dynamical system is simultaneously constrained by independent resolved relative astrometry and \emph{both} SB2 velocity curves. We do not claim priority for this intersection; the literature and source audit is maintained in \texttt{docs/01\_literature\_gap.md} and should be repeated before submission.

The stronger scientific question does not depend on priority. We ask: \emph{what biases in individual stellar dynamical masses and parallax arise when marginal-resolution Gaia measurements are treated as an ordinary photocentre, and how accurate must the astrometric response model be before those biases are removed?} This question naturally leads to a response-fidelity hierarchy. We compare an ordinary photocentre ($\mzero$), the published one-dimensional equal-width blend response ($\mone$), and a finite-elongation orientation-dependent Gaussian surrogate based on \citet{Penoyre2026} ($\mtwo$), while keeping the binary dynamics and the independent resolved-astrometry/SB2 data model fixed.

The experiment sequence is designed to separate three effects. First, a 720-fit paired response-fidelity grid measures how $\mzero$, $\mone$, and $\mtwo$ propagate into the physical parameters as light fraction, angular scale, and profile elongation vary. Second, a 270-fit control weakens only the independent resolved-astrometry and SB2 precision to determine when Gaia response fidelity begins to control the dynamical masses rather than mainly the light ratio. Third, a 100-seed matched-$\mtwo$ control tests whether small offsets seen in a ten-seed grid are repeatable estimator biases or finite-ensemble fluctuations. The earlier 23\,000-fit full-sky and 11\,040-fit matched-transit-count experiments are retained as supporting tests of scan-schedule dependence, not as the central scientific result.

The Gaia-response fidelity experiments remain synthetic: their purpose is to establish inference behaviour and the measurement-model hierarchy before epoch/image products and the associated calibration information support a direct Gaia likelihood test. Separately, the Newtonian resolved-astrometry+SB2 core is tested on the real binary GJ~765.2/HIP~96656 in Section~\ref{sect:realdata}. That V6a experiment is a real-data implementation check of the physical orbit engine and must not be interpreted as validation of the marginal-resolution Gaia response.
